% Measured height against the two analytic bounds. The horizontal axis is
% logarithmic, which is why all three series come out as straight lines.
%
% The gap that matters is the one between the measured series and the upper
% bound: the bound is correct but it is not tight, because the minimum-node
% trees that attain it are vanishingly unlikely to arise from random keys.
\begin{tikzpicture}
  \begin{axis}[
      avlaxis,
      xmode        = log,
      log basis x  = 10,
      xlabel       = {number of keys $n$},
      ylabel       = {height $h$},
      xmin         = 8, xmax = 1.4e6,
      ymin         = 0, ymax = 30,
      legend pos   = north west,
    ]

    \addplot[avlbound] table[x = n, y = upper] {data/height_vs_n.dat};
    \addlegendentry{upper bound $1.4404\lgg(n+2)-1.3277$}

    \addplot[avlmeasured] table[x = n, y = meanh] {data/height_vs_n.dat};
    \addlegendentry{measured height}

    \addplot[avlfloor] table[x = n, y = lower] {data/height_vs_n.dat};
    \addlegendentry{lower bound $\lgg(n+1)$}

  \end{axis}
\end{tikzpicture}
